% Generated draft: requires mathematical and source review.
\documentclass[11pt]{article}
\usepackage[margin=25mm]{geometry}
\usepackage{fontspec}
\setmainfont{Latin Modern Roman}
\usepackage{amsmath,amssymb,mathtools,unicode-math}
\setmathfont{Latin Modern Math}
\usepackage{enumitem,needspace,xurl,hyperref}
\setlength{\parindent}{0pt}
\setlength{\parskip}{5pt}
\setlength{\emergencystretch}{3em}
\newcommand{\N}{\mathbb{N}}
\newcommand{\Z}{\mathbb{Z}}
\newcommand{\Q}{\mathbb{Q}}
\newcommand{\R}{\mathbb{R}}
\newcommand{\C}{\mathbb{C}}
\newcommand{\F}{\mathbb{F}}
\newcommand{\A}{\mathbb{A}}
\newcommand{\PP}{\mathbb{P}}
\newcommand{\E}{\mathbb{E}}
\newcommand{\eps}{\varepsilon}
\newcommand{\dd}{\,\mathrm{d}}
\begin{document}
\section*{UnsolvedMath ID: 801020}
\textbf{Source problem code:} catalogue entry
\subsection{Definitions and mathematical statement}
For an arithmetic-hierarchy index \(r\), \(\Sigma_r\) conventionally denotes the corresponding class of arithmetical formulas. The notation \(\mathrm{I}\Sigma_r\) denotes arithmetic with induction for \(\Sigma_r\) formulas, and \(\mathrm{B}\Sigma_r\) denotes the corresponding bounding or collection scheme. The underlying arithmetic language, base axioms, and parameter conventions are not specified in the supplied record. Write \(T\vdash S\) for syntactic provability of the statement or scheme \(S\) from \(T\); addition of theories denotes combination of their axioms.

The source uses \(\mathrm{WPHP}(\Sigma_{n+1})\) and \(\mathrm{FRT}^{e}_{k}(\Sigma_{n+1})\) without defining them. These labels suggest a weak pigeonhole principle and a finite Ramsey principle restricted by formula complexity, respectively, but their exact schemes cannot be recovered from this record. In particular, their internal quantifiers, numerical bounds, definability requirements, and background theory remain unspecified. The notation is therefore retained without selecting a formulation.

For the index \(n\) appearing in the source, the two requested provability questions are
\[
\mathrm{I}\Sigma_{n+1}+\mathrm{WPHP}(\Sigma_{n+1})
\ \stackrel{?}{\vdash}\ 
\mathrm{FRT}^{2}_{2}(\Sigma_{n+1})
\]
and
\[
\mathrm{FRT}^{2}_{2}(\Sigma_{n+1})
\ \stackrel{?}{\vdash}\ 
\mathrm{FRT}^{3}_{2}(\Sigma_{n+1}).
\]
The admissible range and intended quantification of \(n\) are not supplied. In the second question, any implicit background arithmetic must also be specified before the question determines a precise implication between theories.

As background, the supplied source reports the following nonprovability assertions as known:
\[
\mathrm{I}\Sigma_n+\mathrm{WPHP}(\Sigma_{n+1})
+\mathrm{FRT}^{e}_{k}(\Sigma_{n+1})
\nvdash \mathrm{B}\Sigma_{n+1}
\qquad(e,k\text{ standard}),
\]
\[
\mathrm{FRT}^{e}_{k}(\Sigma_{n+1})
\nvdash \mathrm{WPHP}(\Sigma_{n+1}).
\]
These are supplied background claims, not independently verified results here. The source does not clarify whether the first assertion concerns each fixed standard pair \((e,k)\) or a theory containing all such instances, nor does it explicitly specify the index restrictions for the second assertion. Consequently, this draft preserves the questions and their qualifications but does not supply a fully determined axiomatization.
\subsection{Short English statement}
The source asks whether induction at level \(\Sigma_{n+1}\), together with the indicated weak pigeonhole principle, proves the indicated finite Ramsey principle for pairs and two colors, and whether that Ramsey principle proves the corresponding principle for triples and two colors. These descriptions follow the apparent notation; exact principle definitions, the range of \(n\), and the implicit base theory require clarification. The registry labels the problem open; this is a supplied status, not an independently verified assessment.
\subsection{Sources}
\begin{itemize}
\item Local UnsolvedMath registry, supplied complete record: permanent ID \(801020\), exact source problem code \texttt{catalogue entry}, title ``Open Problems in Mathematical Logic --- Problem 32,'' in the AMR Open Problem Lists collection. The record attributes the item to Barmpalias, Gao et al., with proposed year 2026, and labels it open.
\item Supplied source reference: Barmpalias, Gao et al., \emph{Open Problems in Mathematical Logic} (2026), Problem 32, \url{https://arxiv.org/abs/2608.26628}. The record reports access on 2026-10-16 and extraction by \texttt{arxiv\_latex}; these are supplied provenance statements, not evidence of independent source consultation here.
\item Supplied registry provenance: upstream ID \(801020\), upstream revision \texttt{d371ee917ac9647fbb06cd6fc529deec6f3dda82}; record created and updated at \texttt{2026-10-16T14:00:00Z}. No separate citations supporting the reported background nonprovability results are supplied.
\end{itemize}
\end{document}
